% Part: history
% Chapter: biographies
% Section: gerhard-gentzen

\documentclass[../../../include/open-logic-section]{subfiles}

\begin{document}

\olfileid{his}{bio}{gen}

\olsection{గెర్హార్డ్ గెంట్సెన్}

\olphoto{gentzen-gerhard}{గెర్హార్డ్ గెంట్సెన్}

గెర్హార్డ్ గెంట్సెన్ ప్రధానంగా నిర్మాణాత్మక నిరూపణ సిద్ధాంతాన్ని సృష్టించిన వ్యక్తిగా, ప్రత్యేకించి సహజ నిగమనం, సీక్వెంట్ కలనశాస్త్ర \tetoken{నిగమన}{!!{derivation}} వ్యవస్థలను రూపొందించిన వ్యక్తిగా ప్రసిద్ధి చెందారు. ఆయన 1909 నవంబరు 24న జర్మనీలోని గ్రైఫ్స్‌వాల్డ్‌లో జన్మించారు. సన్నాహక పాఠశాలకు వెళ్లే ముందు గెర్హార్డ్ మూడు సంవత్సరాలు ఇంట్లోనే చదువుకున్నారు; ఆ పాఠశాలలో చదువు పరంగా సహవిద్యార్థులలో చాలామంది కంటే వెనుకబడి ఉన్నారు. అయినా ఆయన ప్రతిభావంతమైన విద్యార్థి, గణితంలో మంచి సామర్థ్యం కనబరిచారు. ఆయన ఆసక్తులు విస్తృతం; ఉదాహరణకు, తల్లి కోసం కవితలు, పాఠశాల నాటకశాల కోసం నాటకాలు కూడా రాశారు.

గెంట్సెన్ విశ్వవిద్యాలయ చదువును గ్రైఫ్స్‌వాల్డ్ విశ్వవిద్యాలయంలో మొదలుపెట్టి, తరువాత గ్యోటింగెన్, మ్యూనిక్, బెర్లిన్ నగరాల్లో చదివారు. హెర్మాన్ వెయిల్ పర్యవేక్షణలో 1933లో గ్యోటింగెన్ విశ్వవిద్యాలయం నుంచి డాక్టరేట్ పొందారు. (ఆయన పనిలో ఎక్కువ భాగాన్ని పాల్ బెర్నాయ్స్ పర్యవేక్షించారు; కానీ నాజీలు ఆయనను విశ్వవిద్యాలయం నుంచి తొలగించారు.) 1934లో గెంట్సెన్ డేవిడ్ హిల్బర్ట్‌కు సహాయకునిగా పని మొదలుపెట్టారు. అదే సంవత్సరంలో \emph{Untersuchungen \"{u}ber das logische Schlie\ss en I--II
  [Investigations Into Logical Deduction I--II]} అనే తన వ్యాసాలలో సీక్వెంట్ కలనశాస్త్రం, సహజ నిగమనం అనే \tetoken{నిగమన}{!!{derivation}} వ్యవస్థలను రూపొందించారు. 1936లో పియానో స్వయంసిద్ధాల అవిరోధత్వాన్ని నిరూపించారు.

నాజీలతో గెంట్సెన్ సంబంధం సంక్లిష్టమైనది. ఆయన మార్గదర్శి బెర్నాయ్స్ జర్మనీని విడిచి వెళ్లాల్సి వచ్చిన సమయంలోనే, గెంట్సెన్ నాజీ అర్ధసైనిక సంస్థ ఎస్‌ఏలోని విశ్వవిద్యాలయ విభాగంలో చేరారు. అనేకమంది జర్మన్లలాగే ఆయన కూడా నాజీ పార్టీలో సభ్యుడు. యుద్ధ సమయంలో వాయు గూఢచార విభాగానికి దూరసంప్రేషణ అధికారిగా పనిచేశారు. అయితే 1942లో నాడీ సంబంధిత కుంగుబాటు కారణంగా విధుల నుంచి విడుదలయ్యారు. గెంట్సెన్ నిజమైన విధేయత నాజీ పార్టీ పట్లే ఉండేదా, లేక విద్యా రంగంలో పురోగతి కోసం పార్టీలో చేరారా అన్నది స్పష్టంగా తెలియదు.

1943లో గెంట్సెన్‌కు ప్రాగ్ జర్మన్ విశ్వవిద్యాలయం గణిత సంస్థలో విద్యాసంబంధ పదవి లభించింది; ఆయన దాన్ని అంగీకరించారు. అయితే 1945లో ప్రాగ్ పౌరులు జర్మన్ ఆక్రమణకు వ్యతిరేకంగా తిరుగుబాటు చేశారు. సోవియట్ దళాలు నగరానికి వచ్చి విశ్వవిద్యాలయ ఆచార్యులందరినీ అరెస్టు చేశాయి. నాజీ సంస్థల్లో సభ్యత్వం కారణంగా గెంట్సెన్‌ను బలవంతపు శ్రమ శిబిరానికి తరలించారు. 1945 ఆగస్టు 4న, 35 ఏళ్ల వయసులో, తన నిర్బంధ గదిలో పోషకాహార లోపంతో మరణించారు.

\begin{reading}
గెంట్సెన్ సమగ్ర జీవిత చరిత్రకు \citet{Menzler-Trott2007} చూడండి. నాజీ పాలనలో గణితశాస్త్రజ్ఞుల గురించి, గెంట్సెన్ జీవితంపై సంక్షిప్త గమనికతో కూడిన రచనకు \citet{Segal2014} చూడండి. తార్కిక నిగమనం గురించి గెంట్సెన్ రాసిన వ్యాసాలు మూల జర్మన్ భాషలో అందుబాటులో ఉన్నాయి \citep{Gentzen1935a,Gentzen1935b}. ఆయన వ్యాసాల ఆంగ్ల అనువాదాలను \citet{Gentzen1969} ఒకే సంపుటంగా సంకలనం చేశారు; అందులో జీవిత చరిత్ర సంక్షిప్త రూపమూ ఉంది.
\end{reading}

\end{document}
